
En d'%
autres termes, @@M0@@ est la @@M1@@-@@M2@@-algèbre engendrée par les générateurs @@M3@@ soumis aux relations (cf. V 4.6)~:
\begin{enumerate}[label=\textup{\roman*)}]
\item @@M4@@ pour @@M5@@ ;
\item @@M6@@.
\end{enumerate}

On commence par le cas où tous les @@M7@@ sont égaux à @@M8@@ ; le théorème s'énonce alors~:

\textbf{Corollaire 4.7.} \uline{Soient} @@M9@@ \uline{un schéma quasi}-\uline{compact}, @@M10@@ \uline{un} @@M11@@\uline{-Module localement libre de rang} @@M12@@, @@M13@@ \uline{le fibré en drapeaux de} @@M14@@, @@M15@@ \uline{les classes des faisceaux inversibles} @@M16@@ (\uline{notation de} 4.5). \uline{Alors} @@M17@@ \uline{est la} @@M18@@-\uline{algèbre engendrée par les} @@M19@@, \uline{soumis aux relations}~:
@@M20@@

On procède par récurrence sur @@M21@@. Pour @@M22@@, le théorème est trivial. Supposons qu'il soit montré pour @@M23@@ ; on considère alors @@M24@@, et sur @@M25@@ la suite exacte~:
@@M26@@
@@M27@@ étant le morphisme structural de @@M28@@. D'après le lemme 4.2, @@M29@@ s'identifie au fibré en drapeaux usuel @@M30@@ sur @@M31@@, les faisceaux canoniques se correspondant dans cette identification. Le diagramme~:
@@M32@@
montre que @@M33@@ se factorise par @@M34@@. D'après l'hypothèse de récurrence appliquée à @@M35@@, et le corollaire 4.4, @@M36@@ admet pour base sur @@M37@@ les éléments de la forme~:
@@M38@@
par ailleurs, @@M39@@ admet pour base sur @@M40@@ les éléments @@M41@@ ; donc @@M42@@ admet pour base sur @@M43@@ les éléments~:
